\documentclass[11pt]{article}
\usepackage{fontspec}
\usepackage[french]{babel}
\usepackage{amsmath}
\usepackage{amssymb}
\usepackage{amsthm}
\usepackage{mathrsfs}
\usepackage{tikz-cd}
\usepackage{enumitem}

\newcommand{\printedpage}[1]{\marginpar{\footnotesize p.~#1}}
\newcommand{\msp}[1]{\marginpar{\footnotesize [#1]}}
\newcommand{\Hom}{\operatorname{Hom}}
\newcommand{\card}{\operatorname{card}}

\begin{document}

\printedpage{199}
% head: A. Grothendieck : Dérivateurs, Ch. XVIII, Édition des Universités de Montpellier II et Paris VII
% foot: Version 17 décembre 2025

\noindent\textbf{Corollaire.} \textit{Soit $\pi$ un cardinal infini. Pour que la sous-catégorie fermée $\underline{S}'$ de $\mathcal{M}$ (formée de $k$-algèbres commutatives $A$ telles que $\Hom_{\mathcal{M}}(A,k)=\emptyset$) soit $\pi$-accessible, il faut et il suffit que pour tout $k$-algèbre $A$ (de cardinal ${}>\pi$), telle que $\Hom_{\mathcal{M}}(A,k)=\emptyset$, existe une sous-algèbre $A_i$} de cardinal $\leqslant\pi$, \textit{telle que $\Hom_{\mathcal{M}}(A_i,k)=\emptyset$.} % [sic] tout k-algèbre -> toute k-algèbre, « algèbre » est féminin

\noindent(On applique le lemme, en prenant pour $\alpha$ le premier ordinal infini.)

\noindent\textbf{Question.} \textit{Cela est-il vérifié, en prenant simplement $\pi\geqslant\card k$~?} La question de l'existence d'un cardinal $\pi$, ayant la propriété indiquée plus haut, me paraît délicate, et j'ignore la réponse. Si on veut construire un contre-exemple pour tout $\pi$, les développements précédents suggèrent de prendre pour $A$ un produit infini d'extensions non triviales de $k$, ou un anneau $k_{\Phi}$ comme dans la démonstration de la proposition 4 a).

\end{document}
